\section{Tổng quan các hướng tiếp cận hỏi--đáp y khoa và xử lý dữ liệu lâm sàng}

Hỏi--đáp trong lĩnh vực y sinh học và hỗ trợ ra quyết định lâm sàng (Clinical Decision Support Systems -- CDSS) là một trong những bài toán phức tạp và đòi hỏi tính chuẩn xác cao nhất của Trí tuệ Nhân tạo hiện đại. Khác với các hệ thống hỏi--đáp mở thông thường, dữ liệu đầu vào trong thực hành y tế lâm sàng mang tính chất không đồng nhất sâu sắc (\emph{Heterogeneous Clinical Data}), bao gồm: các thuộc tính nhân khẩu và tiền sử bệnh nhân ($\mathcal{P}$), văn bản ghi chép diễn tiến khám lâm sàng tự do ($\mathcal{M}_{\mathrm{clin}}$), bản tường trình kết quả chẩn đoán hình ảnh X-quang/CT/MRI dưới dạng văn bản bán cấu trúc ($\mathcal{M}_{\mathrm{img}}$), và bảng kết quả xét nghiệm sinh hóa, huyết học ($\mathcal{L}$).

Lịch sử phát triển của các hệ thống hỏi--đáp y khoa và hỗ trợ quyết định lâm sàng trải qua bốn giai đoạn tiến hóa chính:
\begin{enumerate}[label=(\arabic*)]
    \item \textbf{Hệ chuyên gia dựa trên luật và bản thể học truyền thống (Rule-based \& Ontology):} Tiêu biểu như hệ thống MYCIN hoặc các mô hình luật quyết định lâm sàng \cite{rule_based_1, rule_based_2}. Ưu điểm tuyệt đối là tính minh bạch, suy diễn logic xác định và tuân thủ các quy tắc an toàn y khoa nghiêm ngặt. Tuy nhiên, hệ thống này cực kỳ cứng nhắc, chi phí duy trì cơ sở luật tăng theo hàm mũ và hoàn toàn bất lực trước dữ liệu văn bản tự do trong hồ sơ bệnh án điện tử (EHR).
    \item \textbf{Mô hình ngôn ngữ lớn y tế độc lập (Standalone Medical LLMs):} Sự phát triển vượt bậc của kiến trúc Transformer \cite{vaswani2017} và các mô hình nền tảng như Med-PaLM \cite{singhal2023towards}, LLaMA \cite{llama3_2024} đã tạo ra bước ngoặt về năng lực hiểu và trả lời câu hỏi y tế. Tuy nhiên, việc lưu trữ tri thức ngầm định trong các trọng số tham số cố định bộc lộ ba điểm yếu cốt tử: (i) Nguy cơ sinh ảo giác (\emph{hallucination}) đưa ra các chẩn đoán sai lệch nghiêm trọng; (ii) Tri thức bị đóng băng, không cập nhật kịp thời các phác đồ điều trị mới; và (iii) Thiếu cơ chế minh chứng và giải thích trích dẫn cụ thể để bác sĩ đối soát.
    \item \textbf{Kiến trúc Truy xuất tăng cường tạo sinh (Vanilla Medical RAG):} Sự kết hợp giữa bộ truy xuất dữ liệu ngoài và mô hình ngôn ngữ tạo sinh \cite{lewis2020, karpukhin2020} được xem là giải pháp nền tảng nhằm khắc phục ảo giác. Bằng cách ràng buộc câu trả lời vào các đoạn văn bản tài liệu hướng dẫn điều trị đã kiểm chứng, hệ thống có thể cập nhật tri thức động và cung cấp trích dẫn nguồn cho người dùng. Mặc dù vậy, RAG truyền thống phụ thuộc mù quáng vào độ tương đồng cosine bề mặt, dễ kéo về các đoạn văn bản rác (\emph{noisy snippets}) và không nắm bắt được quan hệ cấu trúc giải phẫu.
    \item \textbf{Kiến trúc RAG đa tác nhân kết hợp Đồ thị tri thức và Lượng giá bất định (Modular Graph-guided \& Uncertainty-aware RAG):} Thế hệ nghiên cứu mới nhất (2024--2026) tích hợp đồ thị tri thức y sinh (Ontology Graph), cơ chế kiểm chứng chuỗi suy luận cấp bước (Process Reward Models) và lượng giá độ bất định lâm sàng (Uncertainty Quantification) để chủ động kiểm soát rủi ro và từ chối trả lời (\emph{Selective Abstention}) khi dữ liệu ca bệnh không đủ tin cậy \cite{knowing_when_abstain_2026, knowguard2024, dong2025discuss}.
\end{enumerate}

Để làm rõ phương pháp luận khoa học theo đúng tinh thần chỉ đạo của Giảng viên hướng dẫn: \emph{ở mỗi công đoạn kỹ thuật, việc lập bảng tiêu chí so sánh phải làm nổi bật được các lỗ hổng nghiên cứu (Research Gaps) cụ thể, từ đó chỉ rõ giải pháp đề xuất của đề tài nhằm lấp đầy các khoảng trống đó}. Toàn bộ chương này được phân tích thành \textbf{bốn công đoạn kỹ thuật cốt lõi}:
\begin{itemize}
    \item \textbf{Công đoạn 1:} Các Mô hình Ngôn ngữ Nền tảng trong Y tế (Medical Foundation LLMs) và Phương pháp Tinh chỉnh Tham số (Fine-tuning / PEFT);
    \item \textbf{Công đoạn 2:} Biểu diễn Tri thức Bản thể học (Ontology) và Đồ thị Tri thức Lâm sàng (Clinical Knowledge Graph);
    \item \textbf{Công đoạn 3:} Kỹ thuật Truy xuất Tăng cường (Medical RAG), Đồ thị RAG (GraphRAG) và Hợp nhất Ngữ cảnh Đa nguồn (Heterogeneous Clinical Text Fusion);
    \item \textbf{Công đoạn 4:} Trí tuệ Nhân tạo Giải thích được (Clinical XAI), Lượng giá Độ bất định (Uncertainty Quantification) và Khắc phục Ảo giác (Hallucination Mitigation).
\end{itemize}
